The light beam arriving from \textbf{Vega Star} can be considered paraxial,
due to the distance from it to the observer on Earth. The explanation for the
existence of two distinct images of the star is that the optical axis of the
objective is not parallel with the light beam from the star.

The images on the camera plate are symmetrical placed relative to the
principal optical axis.

% FIGURE NEEDED: ray diagram (Fig. 12) -- inclined parallel beam from Vega
% entering the objective PQ (aperture 2r, optical centre O, point C above O),
% plane mirror MN of length f/2 lying along the optical axis before the
% photographic plate at distance f; direct rays converge to image Sigma_1 at
% r/2 below the axis, rays reflected by the mirror converge to Sigma_2
% symmetrically above the axis (2014 short-problems solutions, p. 23).

Each of the point images of the Vega Star $\Sigma_1$ and $\Sigma_2$ didn't
concentrate the same light fluxes. In the down below figure it can be seen
the sections of the lens which correspond to each image. The sector APBC is
passed by the light which concentrates in the image $\Sigma_2$ and the light
passing the sector ACBQ concentrates into the point image $\Sigma_1$. See the
picture in figure 13.

% FIGURE NEEDED: lens aperture circle (Fig. 13 b) of radius r centred at O,
% horizontal chord AB through C dividing it into the small upper segment S_2
% and the large lower region S_1; angles 60 deg (COB at O) and 30 deg (at B)
% marked, P and Q at top and bottom of the vertical diameter
% (2014 short-problems solutions, p. 24).

The ratio between the light fluxes concentrated into the two image points
will directly depend on the ratio of the two sectors areas.

From the geometry of the figure 2 results:
\begin{align*}
&MN = OM;\quad N\Sigma_1 = OC = \frac{r}{2};\\
&\angle(CBO) = 30^\circ;\quad \angle(BOC) = 60^\circ;\quad
 \angle(AOB) = 120^\circ;
\end{align*}
\[
\frac{S_1}{S_2} = \frac{8\pi + 3\sqrt{3}}{4\pi - 3\sqrt{3}} \approx 4.
\]

Using the Pogson formula:
\[
\log\frac{E_1}{E_2}
 = \log\frac{\dfrac{\sigma T_{\mathrm{V}}^4 \cdot 4\pi R_{\mathrm{V}}^2}
                   {4\pi d_{\mathrm{PV}}^2} \cdot S_1}
            {\dfrac{\sigma T_{\mathrm{V}}^4 \cdot 4\pi R_{\mathrm{V}}^2}
                   {4\pi d_{\mathrm{PV}}^2} \cdot S_2}
 = -0.4(m_1 - m_2);
\]
\[
\log\frac{S_1}{S_2} = -0.4(m_1 - m_2);
\]
\[
m_2 - m_1 = 1.5^{\mathrm{m}}.
\]
